\section{问题分析}\label{sec:problem-analysis}

本题的核心是在观测不完备的条件下，协调定位精度、搜索完整性与任务耗时。示向度误差使单次检测只能确定源所在的角域，有限接收范围又限制了后续检测位置；在多源情形下，机器狗还须兼顾未知源的发现与已知源的清除。因此，本文依次分析定位区域的几何性质、单源主动选点及多源任务调度，并以实际反馈维持各阶段之间的信息一致性。

\subsection{问题一的分析}

问题一首先需要把测向误差转化为源位置约束。每次示向度对应一个由两条边界射线限定的角域，多次观测的交集构成定位区域。由于不同测点的几何关系可能使交集为空或无界，计算直径前应先检验区域的可行性与有界性；在非空有界情形下，可利用凸多边形的顶点求最远点对，避免遍历区域内的连续点集。

直径计算与覆盖检验需要分别处理。最远点对距离只约束点与点之间的间隔，不能直接保证所有候选位置均落在以该线段为直径的圆内。因此，应另行检验各顶点到圆心的距离，并构造反例说明两者的区别。这一区分也为后续清除提供依据：可靠清除应建立在覆盖整个可行区域的距离保证上，不能仅凭区域直径判断。

\subsection{问题二的分析}

问题二的难点在于第二次观测尚未发生，真实源的位置和接收半径均未知。直接沿首测方向接近可能保留较长的定位区域，过大的横向移动又可能导致失去信号。为此，先将目标区域、首测角域与接收半径上界相交，形成包含真实源的候选区域，再利用接收半径下界筛选能够对全部候选源位置保证接收的检测点。

在该可行范围内，依据定位区域的长轴构造具有不同交会角度的有限探点，通过名义观测估计区域收缩，并同时计入移动代价。名义源点只用于比较动作，不替代实际观测；执行检测后，仍须根据真实反馈更新定位区域，并检查是否达到可靠清除尺度。由此将选点的可行性保证与选点效率的实验评价分开。

\subsection{问题三的分析}

问题三同时包含未知源的发现、已知源的定位清除以及任务完成判定。仅清除当前已经发现的源不能保证全清，而逐站重复扫描全部频道又会增加行程与检测费用。因此，先以接收半径下界构造全域发现覆盖，再将剩余覆盖站与已知源组成两类任务，在保证不遗漏源的前提下联合安排访问顺序。

由于新观测会改变定位区域和任务集合，采用滚动规划：每轮暂时冻结任务代表位置及代理费用，以状态压缩和有预算的 A* 搜索比较任务顺序，仅执行首项任务后重新规划。源处理过程中调用第二问的主动定位模块，并利用共享测量与有严格依据的静默推断减少重复检测；定位不足时以有限光学网格完成局部搜索。最后分别检验全清比例、动作计费用时及程序计算开销，通过同场景对照和消融分析判断各策略的实际作用。

\subsection{问题四的分析}

定向源只向发射方向两侧各 $90^\circ$ 的半平面辐射，无信号既可能因为距离超限，也可能因为检测点在发射范围之外。这使第三问的三处做法不再成立：无信号点不能再用来推距离；覆盖站只按距离布置，不能保证发现背对着检测点的源；贴着边界朝外发射的源，在场内任何位置都收不到。因此，问题四需要先把发现条件改成对任意发射方向都成立的几何条件，并允许检测点放到场地外侧；再为无信号反馈找到在两类源上都成立的用法，例如成对布置检测点，以及把确定在接收范围内的无信号点和有信号点联合起来限制方向和位置。任务安排沿用滚动规划的框架，但每一步的候选动作更多，多测一次是否划算不容易靠规则判断，可以用离线模拟的数据训练一个评分来辅助决策，同时用几何条件守住不漏源、不误清的底线。

